\clearpage
\VerdierPhysicalPageAnchor{0136}%
\begingroup
\small
\setlength{\emergencystretch}{1.5em}%
\setlength{\abovedisplayskip}{1pt plus 1pt minus 1pt}%
\setlength{\belowdisplayskip}{1pt plus 1pt minus 1pt}%

\VerdierSemanticAnchor{ch2-2-2-1-fl-source-target-maps}
\noindent par $\mathbf{s}$ et $\mathbf{b}\colon
\operatorname{Fl}(\mathcal C(S^{-1}))\rightrightarrows
\operatorname{Ob}(\mathcal C)$ les applications :

\VerdierSemanticAnchor{ch2-2-2-1-source-target-formula}
\[
  m\in\operatorname{Hom}_{\mathcal C(S^{-1})}(X,Y),\qquad
  \mathbf{s}(m)=X,\qquad \mathbf{b}(m)=Y,
\]

\VerdierSemanticAnchor{ch2-2-2-1-identity-map}
\noindent et par $\mathbf{id}\colon\operatorname{Ob}(\mathcal C)
\to\operatorname{Fl}(\mathcal C(S^{-1}))$ l'application
$X\mapsto Q(\mathbf{id}_X)$.
\VerdierSemanticAnchor{ch2-2-2-1-category-diagram}%
On définit ainsi un diagramme
$\mathcal C(S^{-1})=(\operatorname{Ob}(\mathcal C),
\operatorname{Fl}(\mathcal C(S^{-1})),\mathbf{s},\mathbf{b},\mathbf{id})$.
On a de plus un couple d'applications noté abusivement :

\VerdierSemanticAnchor{ch2-2-2-1-q-pair}
\[
  Q=\bigl(\mathbf{id}_{\operatorname{Ob}(\mathcal C)}\colon
  \operatorname{Ob}(\mathcal C)\to\operatorname{Ob}(\mathcal C),\;
  Q\colon\operatorname{Fl}(\mathcal C)\to
  \operatorname{Fl}(\mathcal C(S^{-1}))\bigr)\,.
\]

\VerdierSemanticAnchor{ch2-2-2-2}
\smallskip
\noindent\textbf{Théorème 2.2.2.}\quad
\begingroup\itshape
\VerdierSemanticAnchor{ch2-2-2-2-a}%
a) Il existe une et une seule loi de composition partiellement définie sur
$\operatorname{Fl}(\mathcal C(S^{-1}))$ telle que :

\VerdierSemanticAnchor{ch2-2-2-2-a-1}
\par\noindent\hspace*{2em}1) $\mathcal C(S^{-1})$ soit une catégorie.

\VerdierSemanticAnchor{ch2-2-2-2-a-2}
\par\noindent\hspace*{2em}2) $Q(s)$ soit un isomorphisme et $Q(s)^{-1}$ son
inverse $(s\in S)$.

\VerdierSemanticAnchor{ch2-2-2-2-a-3}
\par\noindent\hspace*{2em}3) Pour tout diagramme
$D=X\xleftarrow{s}X'\xrightarrow{m}Y$, $s\in S$, on ait
$\overline D=Q(m)Q(s)^{-1}$
$(\overline D\in\operatorname{Hom}_{\mathcal C(S^{-1})}(X,Y)$ est la classe
de $D$).

\VerdierSemanticAnchor{ch2-2-2-2-a-4}
\par\noindent\hspace*{2em}4) $Q$ soit un foncteur.

\VerdierSemanticAnchor{ch2-2-2-2-b}
\par\noindent b) Tout morphisme de $\mathcal C(S^{-1})$ peut s'écrire comme
un composé $Q(t)^{-1}Q(m)$, $m\in\operatorname{Fl}(\mathcal C)$, $t\in S$,
(resp. $Q(p)Q(s)^{-1}$, $p\in\operatorname{Fl}(\mathcal C)$, $s\in S$).

\VerdierSemanticAnchor{ch2-2-2-2-c}
\par\noindent c) Le foncteur $Q$ commute aux limites projectives finies et
aux limites inductives finies.

\VerdierSemanticAnchor{ch2-2-2-2-d}
\par\noindent d) Pour toute catégorie $\mathcal A$, la composition avec le
foncteur $Q$ définit un foncteur :

\VerdierSemanticAnchor{ch2-2-2-2-hom-q-a}
\[
  \mathcal{H}\!om(Q,\mathcal A)\colon
  \mathcal{H}\!om(\mathcal C(S^{-1}),\mathcal A)
  \longrightarrow \mathcal{H}\!om(\mathcal C,\mathcal A)
\]

\noindent\textup{($\mathcal{H}\!om$ désigne la catégorie des foncteurs).}
\VerdierSemanticAnchor{ch2-2-2-2-universal-image}%
Le foncteur $\mathcal{H}\!om(Q,\mathcal A)$ est injectif sur les objets,
pleinement fidèle, et l'image de
$\mathcal{H}\!om(\mathcal C(S^{-1}),\mathcal A)$ est la sous-catégorie
strictement pleine de $\mathcal{H}\!om(\mathcal C,\mathcal A)$ définie par les
foncteurs de $\mathcal C$ dans $\mathcal A$ qui transforment les morphismes
de $S$ en isomorphismes.
\par\endgroup

\VerdierSemanticAnchor{ch2-2-2-2-proof-omitted}
La démonstration de ce théorème est laissée au lecteur, qui pourra consulter
[16].
\VerdierSemanticAnchor{ch2-2-2-2-composition-setup}%
Nous nous bornerons à décrire la composition des morphismes dans
$\mathcal C(S^{-1})$. Soient
$\overline D_1\in\operatorname{Hom}_{\mathcal C(S^{-1})}(X,Y)$ et
$\overline D_2\in\operatorname{Hom}_{\mathcal C(S^{-1})}(Y,Z)$ deux
morphismes composables de $\mathcal C(S^{-1})$.

\VerdierSemanticAnchor{ch2-2-2-2-roof-pair}
Soient $D_1=X\xleftarrow{s}X'\xrightarrow{m}Y$ et
$D_2=Y\xleftarrow{t}Y'\xrightarrow{p}Z$ deux diagrammes de $\mathcal C$
dans les classes $\overline D_1$ et $\overline D_2$.
\VerdierSemanticAnchor{ch2-2-2-2-sm2-completion-input}%
En utilisant SM2), complétons le diagramme :

\begin{center}
  \begin{tikzpicture}[>=stealth,line width=.38pt,x=1.42cm,y=.90cm,
    scale=.92,transform shape]
    \node (Xp) at (-1.20,1.10) {$X'$};
    \node (Yp) at ( 1.20,1.10) {$Y'$};
    \node (Y)  at ( 0,-.70) {$Y$};
    \draw[->] (Xp) -- node[left=2pt] {$m$} (Y);
    \draw[->] (Yp) -- node[right=2pt] {$t$} (Y);
    \node at (3.00,.25) {$t\in S$};
  \end{tikzpicture}
\end{center}

\VerdierSemanticAnchor{ch2-2-2-2-p137-boundary}
\endgroup
